\documentclass[11pt]{amsart}
\usepackage[margin=1in]{geometry}
\usepackage{amssymb,amsmath,booktabs,longtable}
\usepackage[colorlinks=true,linkcolor=blue,citecolor=blue,urlcolor=blue]{hyperref}
\newtheorem{theorem}{Theorem}[section]
\newtheorem{lemma}[theorem]{Lemma}
\newtheorem{proposition}[theorem]{Proposition}
\newtheorem{corollary}[theorem]{Corollary}
\theoremstyle{definition}
\newtheorem{definition}[theorem]{Definition}
\theoremstyle{remark}
\newtheorem{remark}[theorem]{Remark}
\newcommand{\cA}{\mathsf{A}}\newcommand{\cB}{\mathsf{B}}\newcommand{\cC}{\mathsf{C}}
\newcommand{\Pz}{P_0}
\newcommand{\tri}[3]{\langle #1,\,#2,\,#3\rangle}
\newcommand{\Bset}{\mathcal B}

\title{The three-colour Rado numbers of $ax+3y=3z$}
\author{Puyi Sun}
\date{Version 0.1.1, 20 October 2026 (first public version 18 October 2026). Computer-assisted proof candidate; not peer reviewed.}

\begin{document}
\begin{abstract}
For every integer $a\ge4$ with $3\nmid a$ we prove that $R_3(ax+3y=3z)=a^3+a^2+7a+1$, the case $b=3$ of a
conjecture of Ahmed, Zaman and Bright. The lower bound is their Theorem~1.3; this document proves the upper bound.
The proof combines elementary arguments with $124$ propositional unsatisfiability certificates in LRAT format,
each checked by two independent proof checkers, one of them formally verified. A fixed finite system of constraints
whose vertices are integer polynomials in $k=\lceil a/3\rceil$ handles all large $a$ simultaneously.
\end{abstract}
\maketitle

\tableofcontents

%=====================================================================
\section{Statement}

For a linear equation $\mathcal E$ in three variables, the $3$-colour Rado number $R_3(\mathcal E)$ is the least
positive integer $N$ such that every colouring of $[1,N]=\{1,\dots,N\}$ with three colours contains positive integers
$x,y,z\in[1,N]$ of the same colour satisfying $\mathcal E$. The integers $x,y,z$ need not be distinct.

\begin{theorem}[Ahmed--Zaman--Bright {\cite[Thm.~1.3]{AZB}}]\label{thm:azb}
For coprime integers $a>b\ge3$ with $a^2+a+b>b^2+ba$, $R_3(ax+by=bz)\ge a^3+a^2+(2b+1)a+1$.
\end{theorem}

For $b=3$ the hypotheses mean $a\ge4$ and $3\nmid a$. Ahmed, Zaman and Bright conjectured equality
\cite[Conj.~1.1]{AZB}.

\begin{theorem}\label{thm:main}
For every integer $a\ge4$ with $3\nmid a$,
\[R_3(ax+3y=3z)=a^3+a^2+7a+1 .\]
\end{theorem}

By Theorem~\ref{thm:azb} it remains to prove the upper bound:

\begin{theorem}[upper bound]\label{thm:upper}
Let $a\ge4$, $3\nmid a$, and $L=a^3+a^2+7a+1$. Every $3$-colouring of $[1,L]$ contains a monochromatic solution of
$ax+3y=3z$.
\end{theorem}

\subsection*{Related work}
Since $ax+by=bz$ is equivalent to $b(z-y)=ax$, the numbers $R_3(ax+by=bz)$ are the numbers $R_3(A(x-y)=Bz)$
studied by Chang, De Loera and Wesley \cite{CDW} with $(A,B)=(b,a)$. They proved $R_3(A(x-y)=Bz)=A^3$ for
$\gcd(A,B)=1$ and $A\ge B+2$, and $R_3(A(x-y)=(A-1)z)=A^3+(A-1)^2$ for $A\ge3$; in the present notation these cover
$a\le b-1$ with $\gcd(a,b)=1$ and $b\ge3$, whereas Theorem~\ref{thm:main} concerns $a>b=3$. Their table of exact values contains
$R_3(3(x-y)=Bz)$ for $B\le15$, which agrees with $a^3+a^2+7a+1$ for $4\le a\le14$, $3\nmid a$; Ahmed, Zaman and
Bright computed the values for $a\le30$ \cite{AZB}. The method of Lemma~\ref{lem:W}, a propositional formula
whose variables are indexed by integer polynomials in a parameter, is of the same kind as the symbolic encoding of
\cite[Lemma~4.5]{CDW}; see also the symbolic sets of \cite{AZB}.

\subsection*{Structure of the proof}
\begin{enumerate}
\item $4\le a\le 26$: sixteen certificates (Section~\ref{sec:small}).
\item $a\ge13$, Lemma~\ref{lem:W} (computer-assisted): the colours of $3,6,\dots,24$ form one of $53$ patterns.
\item $a\ge23$, Lemma~\ref{lem:r76} (computer-assisted, $106$ certificates independent of $a$) and
Lemma~\ref{lem:e4}: the pattern is $\Pz$.
\item $a\ge7$, Section~\ref{sec:p0} (no computer): the pattern $\Pz$ forces the colours of all multiples of $3$ up to
$3(a^2+a+1)$ and then $L$ admits no colour.
\end{enumerate}
The thresholds $13,23,7$ are all at most $26$, so (1) and (2)--(4) together cover all $a\ge4$ with $3\nmid a$.

%=====================================================================
\section{Notation and the three families of constraints}\label{sec:setup}

Throughout, $a\ge4$ and $3\nmid a$. Put
\[H=a^2+a+1,\qquad M=H-3=a^2+a-2,\qquad F=H+5=a^2+a+6,\]
\[L=1+a(H+6)=a^3+a^2+7a+1 .\]

\subsection{Reduction of the equation}
If $ax+3y=3z$ then $3\mid ax$, and since $3\nmid a$ we get $3\mid x$. Writing $x=3w$ gives $z=y+aw$.

Let $c\colon[1,L]\to\{\text{three colours}\}$ be a colouring. We only look at two arithmetic progressions:
\[h(t)=c(3t)\quad(1\le t\le H),\qquad f(q)=c(1+aq)\quad(0\le q\le H+6).\]
Note $3H=3a^2+3a+3\le L-1$, $1+aF=L-a\le L-1$ and $1+a(H+6)=L$. Thus $h$ and $f|_{[0,F]}$ only involve integers in
$[1,L-1]$, and $f(H+6)=c(L)$.

\subsection{The constraints}
We use the following solutions of $ax+3y=3z$; in each case all three entries lie in $[1,L]$.

\begin{description}
\item[(BB)] For integers $u\ge1$, $s\ge1$ with $3u\le H$ and $s+au\le H$: $(x,y,z)=(9u,\,3s,\,3(s+au))$. Indeed
$a\cdot9u+3\cdot3s=3\cdot3(s+au)$. Hence $h(3u),h(s),h(s+au)$ are not all equal. We write such a triple of
$h$-indices as $\tri{3u}{s}{s+au}$ and call it a \emph{BB-triple}; its first entry is a multiple of $3$ and its third
entry is the second plus $a\cdot(\text{first}/3)$.
\item[(AB)] For integers $t\ge1$, $q\ge0$ with $t\le H$ and $q+t\le H+6$: $(x,y,z)=(3t,\,1+aq,\,1+a(q+t))$. Indeed
$3at+3(1+aq)=3(1+a(q+t))$. Hence $h(t),f(q),f(q+t)$ are not all equal. In words: \emph{two positions $q<q'$ of $f$
with $f(q)=f(q')$ must not have this common colour equal to $h(q'-q)$.}
\item[(ID)] If $3s=1+aq$ with $1\le s\le H$ and $0\le q\le H+6$, then $h(s)=f(q)$, since both are $c(3s)$.
\end{description}

A BB-triple may have repeated entries (for example $3u=s$); it then constrains only two values of $h$, which is
legitimate because solutions need not consist of distinct integers.

A colouring is \emph{solution-free} on a set $S$ if no monochromatic solution has all three entries in $S$.

\subsection{Relaxations}\label{sec:relax}
Each computational model below is a propositional formula whose variables encode colours of some of the integers
above. Its clauses are of two kinds:
\begin{itemize}
\item \emph{base clauses}: each variable takes exactly one colour, and some of the constraints (BB), (AB), (ID);
\item \emph{case clauses}, described explicitly with each model: a symmetry normalisation, a prescribed prefix, a
first-occurrence condition, or a target condition (possibly built with auxiliary variables that are only required to
imply their defining conditions).
\end{itemize}
Each model is used to rule out one specified class of colourings, namely the solution-free colourings of $[1,L-1]$
(resp.\ $[1,L]$) that satisfy its case conditions (after renaming colours where a normalisation is used). A
colouring in that class satisfies every base clause, because each base clause expresses a genuine solution or the
fact that each integer has one colour; it satisfies the case clauses by assumption, setting each auxiliary
variable to true exactly when its defining condition holds. Hence it induces a satisfying assignment, and if the
model is unsatisfiable the class is empty. That the classes ruled out in this way cover all cases is argued
separately in each lemma.

In the polynomial models of Section~\ref{sec:W} two distinct polynomials in $k$ are treated as distinct variables
even when they take equal values for some $k$. A real colouring assigns equal colours to equal integers, so it still
induces an assignment satisfying all base clauses; the model merely omits the constraints ``these two variables are
equal''.

\subsection{Words, prefixes and $h^\ast$}
Colours may be renamed freely. For a word $w_1w_2\cdots w_n$ over the three colours, its \emph{normalisation} is
obtained by renaming the colours $\cA,\cB,\cC$ in order of first appearance. The \emph{prefix} of a colouring is the
normalisation of $h(1)h(2)\cdots h(8)$. Let
\[\Pz=\cA\cA\cB\cA\cA\cB\cA\cA .\]
Let $\Bset=\{3,6,9,\dots,3a\}$ (the multiples of $3$ in $[3,3a]$) and define $h^\ast(t)=\cB$ for $t\in\Bset$ and
$h^\ast(t)=\cA$ otherwise.

%=====================================================================
\section{Small $a$}\label{sec:small}

\begin{lemma}\label{lem:small}
Theorem~\ref{thm:upper} holds for $a\in\{4,5,7,8,10,11,13,14,16,17,19,20,22,23,25,26\}$.
\end{lemma}
\begin{proof}[Proof (computer-assisted)]
Let $U(a)=\{3t:1\le t\le H\}\cup\{1+aq:0\le q\le H+6\}\subseteq[1,L]$ (its largest element is $L$). The formula
$\mathtt{U}a$ has three Boolean variables per element of $U(a)$, the clauses ``at least one colour'' and ``at most
one colour'' for each element, and, for every ordered triple $(x,y,z)\in U(a)^3$ with $ax+3y=3z$ and every colour,
the clause ``$x,y,z$ do not all have this colour''. In addition it fixes $c(1)=\cA$ and $c(3)\in\{\cA,\cB\}$, which
is no loss of generality after renaming colours. The formula is unsatisfiable for each listed $a$
(Section~\ref{sec:cert}). A colouring of $[1,L]$ without monochromatic solution, after renaming its colours so that
$c(1)=\cA$ and $c(3)\in\{\cA,\cB\}$, would restrict to a satisfying assignment.
\end{proof}

From now on $a\ge27$ (the arguments below need only $a\ge23$).

%=====================================================================
\section{The polynomial window (Lemma W)}\label{sec:W}

\subsection{Parametrisation}
For $a\ge13$ with $3\nmid a$ put $k=\lceil a/3\rceil$ and $\rho=a-3k\in\{-2,-1\}$. Then $a=3k+\rho$ and $k\ge5$. The
integer $k$ is determined by $a$; consequently $a$, $ak=3k^2+\rho k$ and $a(k+1)=3k^2+(3+\rho)k+\rho$ are polynomials
in $k$ with integer coefficients, and so is every position used in this section.

\subsection{The model $\mathcal W_\rho$}
\emph{$f$-vertices.} The positions
\[q=\lambda a+3kD+e,\]
for the triples $(\lambda,D,e)$ listed in Table~\ref{tab:window}, regarded as integer polynomials in $k$ after
substituting $a=3k+\rho$. Equal polynomials are identified.

\begin{table}[h]\centering
\begin{tabular}{ccc}\toprule
$\lambda$ & $D$ & range of $e$\\\midrule
$0$ & $0$ & $0\le e\le 7$\\ $0$ & $1$ & $-7\le e\le 7$\\ $0$ & $2$ & $-7\le e\le 0$\\
$k$ & $1$ & $-1\le e\le 7$\\ $k$ & $2$ & $-7\le e\le 7$\\ $k$ & $3$ & $-7\le e\le 4$\\
$k+1$ & $1$ & $-2\le e\le 7$\\ $k+1$ & $2$ & $-7\le e\le 7$\\ $k+1$ & $3$ & $-7\le e\le 4$\\\bottomrule
\end{tabular}
\caption{The $f$-vertices of the window.}\label{tab:window}
\end{table}

\emph{$h$-vertices.} The constants $1,2,\dots,8$; the heads $3,6,3k,3k+3$; and every difference $q'-q$ of two
$f$-vertices which, as a polynomial, takes values $\ge1$ for every integer $k\ge5$.

\emph{Range.} Every $h$-vertex takes values in $[1,H]$ and every $f$-vertex in $[0,F]$ for every integer $k\ge5$.
(For each quadratic this is decided exactly from its leading coefficient and the integer minimiser.) Hence every
constraint below is an instance of (AB) or (BB) inside $[1,L-1]$.

\emph{Clauses.}
\begin{itemize}
\item exactly one colour per vertex;
\item (AB) for every pair of $f$-vertices $q,q'$ with $q'-q\ge1$ for all $k\ge5$: $h(q'-q),f(q),f(q')$ not all equal;
\item (BB) for every head $3u$ with $u\in\{1,2,k,k+1\}$ and every $h$-vertex $s$ such that $s+au$ (as a polynomial)
is also an $h$-vertex: $h(3u),h(s),h(s+au)$ not all equal;
\item the target: one clause, built with auxiliary variables that are only required to imply their defining
conditions, stating that the colouring of $h(1),\dots,h(8)$ uses three colours, or has $h(3)\ne h(6)$, or has as
normalisation one of the eleven words
\[\begin{gathered}
\cA\cA\cA\cA\cA\cA\cA\cA,\ \cA\cA\cA\cA\cA\cA\cA\cB,\ \cA\cA\cA\cA\cA\cA\cB\cA,\ \cA\cA\cA\cA\cB\cA\cA\cA,\
\cA\cA\cA\cA\cB\cA\cB\cA,\ \cA\cA\cA\cB\cA\cA\cA\cA,\\
\cA\cB\cA\cA\cA\cA\cA\cA,\ \cA\cB\cB\cB\cA\cB\cA\cB,\ \cA\cB\cB\cB\cA\cB\cB\cB,\ \cA\cB\cB\cB\cB\cB\cA\cB,\
\cA\cB\cB\cB\cB\cB\cB\cB.
\end{gathered}\]
\end{itemize}
The colours of $h(1),\dots,h(8)$ are not fixed in the model.

\subsection{The lemma}
Call a word $\cA w_2w_3\cdots w_8$ over $\{\cA,\cB\}$ with $w_3=w_6$ \emph{admissible} if it is not one of the eleven
words above. There are $2^6=64$ words $\cA w_2\cdots w_8$ over $\{\cA,\cB\}$ with $w_3=w_6$, hence $53$ admissible
words; they are listed in Appendix~\ref{app:53}, and $\Pz$ is one of them.

\begin{lemma}[Lemma W]\label{lem:W}
Let $a\ge13$, $3\nmid a$, and let $c$ be a solution-free $3$-colouring of $[1,L-1]$. Then the prefix of $c$ is an
admissible word.
\end{lemma}
\begin{proof}[Proof (computer-assisted)]
The prefix of a colouring that uses at most two colours on $h(1..8)$ and has $h(3)=h(6)$ is a word
$\cA w_2\cdots w_8$ over $\{\cA,\cB\}$ with $w_3=w_6$. So if the prefix is not admissible, the target clause holds.
For $\rho=-2$ and $\rho=-1$ the formula $\mathcal W_\rho$ is unsatisfiable (Section~\ref{sec:cert}). Since $c$
induces a satisfying assignment of $\mathcal W_\rho$ whenever the target holds (Section~\ref{sec:relax}), the target
fails, and the prefix is admissible.
\end{proof}

%=====================================================================
\section{From $53$ prefixes to $\Pz$}\label{sec:53}

\subsection{The chain lemma}
\begin{lemma}[chain]\label{lem:chain}
Let $m\ge1$ and let $g\colon[1,m]\to\{X,Y\}$ be a two-colouring with no monochromatic BB-triple inside $[1,m]$.
Suppose $g(3)=X$, $k\ge2$, $k(a+3)+a\le m$, and $g(3j)=X$ for $1\le j<k$. Then $g(3k)=X$.
\end{lemma}
\begin{proof}
Suppose $g(3k)=Y$ and put $N=k(a+3)=3k+ak$. All indices below are at most $N+a\le m$. The following are BB-triples:
\begin{center}\begin{tabular}{lll}\toprule
triple & $u$ & check\\\midrule
$\tri{3k}{3k}{N}$ & $k$ & $3k+ak=N$\\
$\tri{3(k-1)}{3k+a}{N}$ & $k-1$ & $3k+a+a(k-1)=N$\\
$\tri{3}{N}{N+a}$ & $1$ & $N+a=N+a$\\
$\tri{3k}{3k+a}{N+a}$ & $k$ & $3k+a+ak=N+a$\\\bottomrule
\end{tabular}\end{center}
From the first triple, $g(N)=X$. In the second, $g(3(k-1))=X$ (as $k-1\ge1$) and $g(N)=X$, so $g(3k+a)=Y$. In the
third, $g(3)=g(N)=X$, so $g(N+a)=Y$. Now the fourth triple is monochromatic of colour $Y$, a contradiction.
\end{proof}

\subsection{Two colours force $\Pz$}
\begin{lemma}[E4]\label{lem:e4}
Let $a\ge23$, $J=\lceil(2a+8)/3\rceil$ and $m_0=J(a+3)+a$. Let $c$ be a solution-free colouring of $[1,L-1]$ such that
$h$ takes only two colours on $[1,m_0]$. Then the prefix of $c$ is $\Pz$.
\end{lemma}
\begin{proof}
One checks $m_0\le H$ for $a\ge23$ (Section~\ref{sec:ranges}), so $h$ is defined on $[1,m_0]$, and it has no
monochromatic BB-triple there. Let $X=h(3)$ and $Y$ the other colour. Lemma~\ref{lem:chain} with $m=m_0$, applied
for $k=2,3,\dots,J$ in turn, gives $h(3k)=X$ for $1\le k\le J$.

Let $d\in\{1,2,4,5,7,8\}$. Since $3\nmid d$ and $3\nmid a$, there is $u\in\{1,2\}$ with $3\mid au+d$. Put
$v=(au+d)/3$; then $1\le v\le(2a+8)/3\le J$. Now $\tri{3u}{d}{3v}$ is a BB-triple ($d+au=3v\le 3J\le m_0$), and
$h(3u)=h(3v)=X$ because $u,v\le J$. Hence $h(d)=Y$. So $h(1)\cdots h(8)=YYXYYXYY$, whose normalisation is $\Pz$.
\end{proof}

\subsection{The third colour does not appear early}
\begin{lemma}[computer-assisted, independent of $a$]\label{lem:r76}
Let $a\ge23$, let $c$ be a solution-free colouring of $[1,L-1]$ whose prefix $X$ is an admissible word, and let
$\cC$ be the colour not occurring in $X$ (after the normalising renaming). Then $h(t)\ne\cC$ for $9\le t\le m_0$.
\end{lemma}
\begin{proof}[Proof (computer-assisted)]
For each admissible $X$ Appendix~\ref{app:53} gives an integer $K=K_X$ with $10\le K\le 68$. Suppose $\cC$ occurs in
$h(9..m_0)$ and let $t_0$ be its first occurrence. Then either $t_0\le K-1$, or $K\le t_0\le m_0$; these two cases
are exhaustive.

\emph{Early case, $9\le t_0\le K-1$.} Put $F_X=K+79$. The formula $\mathtt{I}_{X}$ has variables for $h(1),\dots,
h(F_X)$ and $f(0),\dots,f(F_X)$ (three colours each, exactly one colour per vertex), the unit clauses $h(1..8)=X$,
the clause ``$h(t)=\cC$ for some $t$ with $9\le t\le K-1$'', and every (AB) constraint
$h(t),f(q),f(q+t)$ not all equal with $1\le t$, $0\le q$, $q+t\le F_X$. All positions satisfy $F_X\le147\le\min(H,F)$
for $a\ge23$, so this is a relaxation, and $\mathtt I_X$ is unsatisfiable.

\emph{Late case, $K\le t_0\le m_0$.} By minimality of $t_0$, $h(t)\in\{\cA,\cB\}$ for $9\le t\le K-1$. Consider the
two rows $f(0),\dots,f(K-1)$ and $f(t_0),\dots,f(t_0+K-1)$. Since $t_0+K-1\le m_0+67\le F$ for $a\ge23$
(Section~\ref{sec:ranges}), all these are values of $f$ on $[0,F]$, and the rows are disjoint because $t_0\ge K$.
The formula $\mathtt R_X$ has one Boolean variable for each $h(t)$, $9\le t\le K-1$ (colour $\cA$ or $\cB$), three
colour variables for each of the $2K$ cells, exactly one colour per cell, and the clauses
\begin{itemize}
\item for each row, each $d\in[1,K-1]$ and each $i\in[0,K-1-d]$: cells $i$ and $i+d$ of the row are not both of
colour $h(d)$ (with $h(d)$ given by $X$ for $d\le8$ and by the Boolean variable for $d\ge9$); this is (AB) with
difference $d$;
\item for each $i\in[0,K-1]$: cell $i$ of the first row and cell $i$ of the second row are not both $\cC$; this is
(AB) with difference $t_0$, because $h(t_0)=\cC$.
\end{itemize}
The formula $\mathtt R_X$ does not depend on $t_0$ or $a$, and it is unsatisfiable.

In both cases the colouring would induce a satisfying assignment, a contradiction.
\end{proof}

\begin{corollary}\label{cor:p0}
Let $a\ge23$ and let $c$ be a solution-free colouring of $[1,L-1]$. Then the prefix of $c$ is $\Pz$.
\end{corollary}
\begin{proof}
By Lemma~\ref{lem:W} the prefix $X$ is admissible, so $h(1..8)$ uses at most the two colours $\cA,\cB$ of $X$. By
Lemma~\ref{lem:r76}, $h$ does not take the third colour on $[9,m_0]$. Hence $h$ uses at most two colours on
$[1,m_0]$, and Lemma~\ref{lem:e4} shows that the prefix is $\Pz$.
\end{proof}

%=====================================================================
\section{The prefix $\Pz$ leads to a contradiction}\label{sec:p0}

Throughout this section $a\ge7$, $c$ is a $3$-colouring of $[1,L]$ without monochromatic solution, and the prefix of
$c$ is $\Pz$. Name the colours so that $h(1..8)=\cA\cA\cB\cA\cA\cB\cA\cA$, and let $\cC$ be the third colour. Thus
\[h(3)=h(6)=\cB,\qquad h(1)=h(2)=h(4)=h(5)=h(7)=h(8)=\cA .\]

\subsection{The third colour is absent below $M$}
\begin{lemma}[nine columns]\label{lem:nine}
If $9\le t\le M$ then $h(t)\ne\cC$.
\end{lemma}
\begin{proof}
Suppose $h(t)=\cC$. Consider the two rows $f(0),\dots,f(8)$ and $f(t),\dots,f(t+8)$; since $t+8\le M+8=F$, these
are defined. Index the nine columns by $i\in[0,8]$, the cells of column $i$ being $f(i)$ and $f(t+i)$.
\begin{enumerate}
\item By (AB) with difference $t$ and $h(t)=\cC$, no column has $\cC$ in both cells.
\item Within a row, two cells at distance $\delta\in[1,8]$ are not both of colour $h(\delta)$. For $\delta\in\{3,6\}$
this colour is $\cB$; for $\delta\in\{1,2,4,5,7,8\}$ it is $\cA$.
\item Split the columns into the three classes $i\equiv0,1,2\pmod 3$. Two cells of one row in the same class are at
distance $3$ or $6$, so each row has at most one $\cB$ in each class. Two cells of one row in different classes are
at a distance in $\{1,2,4,5,7,8\}$, so they are not both $\cA$; hence all cells coloured $\cA$ in one row lie in a
single class.
\item Therefore the cells coloured $\cA$ in the two rows lie in at most two classes, and some class $\gamma$ contains
no $\cA$ in either row. In each row the three cells of class $\gamma$ contain at most one $\cB$ and no $\cA$, hence at
least two $\cC$. So the six cells of class $\gamma$ contain at least four $\cC$, distributed over three columns; some
column has $\cC$ in both cells, contradicting (1).\qedhere
\end{enumerate}
\end{proof}

So $h$ takes only the colours $\cA,\cB$ on $[1,M]$, with $h(3)=\cB$.

\subsection{Two-colour uniqueness}
\begin{theorem}\label{thm:u2c}
Let $a\ge7$ and let $g\colon[1,M]\to\{\cA,\cB\}$ have no monochromatic BB-triple inside $[1,M]$, with $g(3)=\cB$. Then
$g=h^\ast$ on $[1,M]$.
\end{theorem}
\begin{proof}
\emph{Step 1: $g=\cB$ on $\{3,6,\dots,3(a-3)\}$.} Apply Lemma~\ref{lem:chain} with $m=M$ for $k=2,\dots,a-3$; the
largest index needed is $(a-3)(a+3)+a=a^2+a-9\le M$.

\emph{Step 2: $g(3(a-2))=g(3(a-1))=g(3a)=\cB$.} Put
\[d=3(a-2),\ e=3(a-1),\ p=a+3,\ r=a^2-a+3,\ s=a^2+3,\ v=2a-6,\ w=a^2+a-6,\ v_3=a-6 .\]
All these lie in $[1,M]$ for $a\ge7$. The following are BB-triples (the last column verifies $z=y+au$):
\begin{center}\begin{tabular}{llll}\toprule
label & triple & $u$ & check\\\midrule
T1 & $\tri{3}{3}{p}$ & $1$ & $3+a=p$\\
T2 & $\tri{d}{p}{r}$ & $a-2$ & $a+3+a(a-2)=a^2-a+3$\\
T3 & $\tri{3}{r}{s}$ & $1$ & $a^2-a+3+a=a^2+3$\\
T4 & $\tri{e}{3}{r}$ & $a-1$ & $3+a(a-1)=a^2-a+3$\\
T5 & $\tri{e}{p}{s}$ & $a-1$ & $a+3+a(a-1)=a^2+3$\\
T6 & $\tri{3}{v}{d}$ & $1$ & $2a-6+a=3a-6$\\
T7 & $\tri{d}{d}{w}$ & $a-2$ & $3a-6+a(a-2)=a^2+a-6$\\
T8 & $\tri{e}{v}{w}$ & $a-1$ & $2a-6+a(a-1)=a^2+a-6$\\
T9 & $\tri{6}{v_3}{d}$ & $2$ & $a-6+2a=3a-6$\\
T10 & $\tri{3a}{v_3}{w}$ & $a$ & $a-6+a^2=a^2+a-6$\\\bottomrule
\end{tabular}\end{center}
(a) Suppose $g(d)=\cA$. By T1 ($g(3)=\cB$) $g(p)=\cA$; by T2 $g(r)=\cB$; by T3 $g(s)=\cA$; by T4 ($g(3)=g(r)=\cB$)
$g(e)=\cA$; now T5 is monochromatic $\cA$. Hence $g(d)=\cB$.

(b) By T6 ($g(3)=g(d)=\cB$) $g(v)=\cA$; by T7 $g(w)=\cA$; by T8, $g(e)=\cB$.

(c) By T9 ($g(6)=g(d)=\cB$, using $6\in\{3,\dots,3(a-3)\}$) $g(v_3)=\cA$; by T10 ($g(v_3)=g(w)=\cA$), $g(3a)=\cB$.

So $g=\cB$ on $\Bset$.

\emph{Step 3: $g=\cA$ on $[1,M]\setminus\Bset$.} Let $n\in[1,M]\setminus\Bset$. Since $3\nmid a$ there is a unique
$u_0\in\{1,2,3\}$ with $au_0\equiv n\pmod3$.

\emph{Case $n\ge au_0+3$.} For $u\equiv u_0\pmod 3$ with $1\le u\le a$, the numbers $au+3v$ ($1\le v\le a$) are
exactly the numbers $\equiv n\pmod 3$ in the interval $I_u=[au+3,\,au+3a]$. For consecutive such $u$ and $u+3$ the
interval $I_{u+3}$ begins at $a(u+3)+3=(au+3a)+3$, i.e.\ three after the end of $I_u$, so the intervals $I_u$
cover every number $\equiv n\pmod3$ from $au_0+3$ up to $au_{\max}+3a$, where $u_{\max}\ge a-2$ is the largest
admissible $u$. Since $au_{\max}+3a\ge a(a-2)+3a=a^2+a>M\ge n$, we can write $n=au+3v$ with $1\le u,v\le a$. Then
$\tri{3u}{3v}{n}$ is a BB-triple with $3u,3v\in\Bset$, so $g(n)=\cA$.

\emph{Case $n<au_0+3$.} Then $n\le au_0$. If $u_0=3$ then $n\equiv0\pmod3$ and $n\le3a$, i.e.\ $n\in\Bset$, which is
excluded. If $u_0=1$ then $n\le a$ and $n+2a\equiv a+2a\equiv0\pmod3$ with $2a<n+2a\le3a$, so $n+2a\in\Bset$; the
BB-triple $\tri{6}{n}{n+2a}$ gives $g(n)=\cA$. If $u_0=2$ then $n\le2a$ and $n+a\equiv3a\equiv0\pmod3$ with
$a<n+a\le3a$, so $n+a\in\Bset$; the BB-triple $\tri{3}{n}{n+a}$ gives $g(n)=\cA$.
\end{proof}

\begin{corollary}\label{cor:hstar}
$h=h^\ast$ on $[1,M]$.
\end{corollary}
\begin{proof}
By Lemma~\ref{lem:nine} and the prefix, $h$ is a two-colouring of $[1,M]$ by $\cA,\cB$ with $h(3)=\cB$; it has no
monochromatic BB-triple. Apply Theorem~\ref{thm:u2c}.
\end{proof}

\subsection{The top of the range of $h$}
\begin{lemma}[tail]\label{lem:tail}
$h(H-2)=h(H-1)=\cA$ and $h(H)=\cC$. Consequently $h=h^\ast$ on $[1,H-1]$.
\end{lemma}
\begin{proof}
\emph{An auxiliary value.} Let $q_T=11$ if $a\equiv1$ and $q_T=10$ if $a\equiv2\pmod 3$. Then $3\mid1+aq_T$, and
$s_T=(1+aq_T)/3$ satisfies $3a<s_T\le M$ (as $a\ge7$). Hence $h(s_T)=h^\ast(s_T)=\cA$ and, by (ID), $f(q_T)=\cA$.

\emph{$h(H-1)\ne\cB$ and $h(H-2)\ne\cB$.} Note $H-1=a^2+a$ and $H-2=a^2+a-1$.
\begin{itemize}
\item $\tri{3(a-2)}{3a}{H-1}$ is a BB-triple ($u=a-2$: $3a+a(a-2)=a^2+a$) whose first two entries lie in $\Bset$.
\item If $a\equiv1\pmod 3$: $\tri{3a}{a-1}{H-2}$ ($u=a$: $a-1+a^2$); here $a-1\equiv0\pmod 3$ and $3\le a-1\le3a$, so
$a-1\in\Bset$.
\item If $a\equiv2\pmod 3$: $\tri{3(a-1)}{2a-1}{H-2}$ ($u=a-1$: $2a-1+a(a-1)=a^2+a-1$); here $2a-1\equiv0\pmod 3$
and $2a-1\le3a$, so $2a-1\in\Bset$.
\end{itemize}

\emph{$h(t)\ne\cC$ for $t\in\{H-2,H-1\}$.} Suppose $h(t)=\cC$. Consider the rows $f(0),f(3),f(6)$ and
$f(t),f(t+3),f(t+6)$; note $t+6\le H+5=F$. Column $j\in\{0,3,6\}$ consists of $f(j)$ and $f(t+j)$.
\begin{itemize}
\item By (AB) with difference $t$, no column has $\cC$ in both cells.
\item No cell is $\cA$. For the first row, the distances from $q_T$ to $0,3,6$ are $\{11,8,5\}$ or $\{10,7,4\}$; they
are not multiples of $3$, so $h$ has colour $\cA$ at these distances, and since $f(q_T)=\cA$, (AB) forbids $\cA$. For
the second row the distances $t+j-q_T$ lie in $[a^2+a-12,\,a^2+a-4]\subseteq(3a,M]$ (as $a\ge7$), where $h=h^\ast=\cA$,
and the same argument applies.
\item Within a row the distances are $3$ and $6$, where $h=\cB$; so each row has at most one $\cB$.
\end{itemize}
Hence each row has at least two $\cC$, four $\cC$ in three columns, and some column is $\cC$ in both cells: a
contradiction. So $h(H-2)=h(H-1)=\cA$; since $H-2,H-1>3a$ this agrees with $h^\ast$.

\emph{$h(H)\ne\cA$.} $\tri{3a+3}{1}{H}$ is a BB-triple ($u=a+1$: $1+a(a+1)=H$) with $h(3a+3)=h^\ast(3a+3)=\cA$ (as
$3a<3a+3\le M$) and $h(1)=\cA$.

\emph{$h(H)\ne\cB$.} If $a\equiv2\pmod 3$: $\tri{3a}{a+1}{H}$ ($u=a$: $a+1+a^2=H$), and $a+1\in\Bset$. If
$a\equiv1\pmod 3$: $\tri{3(a-1)}{2a+1}{H}$ ($u=a-1$: $2a+1+a(a-1)=H$), and $2a+1\in\Bset$. In both cases the first two
entries lie in $\Bset$.

Hence $h(H)=\cC$.
\end{proof}

\subsection{$L$ has no admissible colour}
\begin{proposition}\label{prop:L}
There is no colour available for $L$.
\end{proposition}
\begin{proof}
Let $q_0\in\{1,2\}$ with $3\mid1+aq_0$ ($q_0=2$ if $a\equiv1$, $q_0=1$ if $a\equiv2\pmod 3$), and let
$Q=\{q_0+3j:0\le j\le a\}$.

\emph{The colours of $f$ on $Q$.} For $q=q_0+3j\in Q$ put $s_j=(1+aq)/3=s_0+aj$. Then $1\le s_j\le s_0+a^2\le
(1+2a)/3+a^2\le M$ (as $a\ge7$), so by (ID) and Corollary~\ref{cor:hstar}, $f(q)=h^\ast(s_j)$. As $3\nmid a$, exactly
one $j\in\{0,1,2\}$ makes $s_j$ divisible by $3$, and for it $s_j\le s_0+2a\le3a$, so $f(q_0+3j)=\cB$. For the other
two $j\in\{0,1,2\}$, and for all $j\ge3$ (where $s_j>3a$), $f(q_0+3j)=\cA$. In particular:
\begin{itemize}
\item one of $f(q_0)$, $f(q_0+3)$ is $\cA$, at a position $p_{\rm low}\le5$;
\item $f(p_{\rm high})=\cA$ for $p_{\rm high}=q_0+3a$ (here $j=a\ge3$), and $p_{\rm high}\le 3a+2<H$.
\end{itemize}

\emph{The rule.} If $f(q)=f(q')=\cA$ with $1\le q'-q\le H-1$, then by (AB) $h(q'-q)\ne\cA$, i.e.\ (by
Lemma~\ref{lem:tail}) $q'-q\in\Bset$. In particular $q'-q$ is a multiple of $3$ at most $3a$.

\emph{(1) Every position $q\in[0,H+6]$ with $f(q)=\cA$ satisfies $q\equiv q_0\pmod 3$.} Otherwise $|q-p_{\rm high}|$
is not a multiple of $3$ and nonzero, so by the rule $|q-p_{\rm high}|\ge H$. If $q>p_{\rm high}$ this gives
$q\ge H+p_{\rm high}\ge H+3a+1>H+6$; if $q<p_{\rm high}$ it gives $q\le p_{\rm high}-H<0$. Both are impossible.

\emph{(2) $f(q)\ne\cA$ for $H\le q\le H+5$.} If $f(q)=\cA$ then $q\equiv q_0\pmod3$ by (1), and $\delta=q-p_{\rm
high}$ is a multiple of $3$ with $\delta\ge H-3a-2=a^2-2a-1>3a$ (as $a\ge6$) and $\delta\le H+5-3a-1<H$. So $\delta\notin\Bset$,
contradicting the rule.

\emph{(3) One of $f(H)$, $f(H+3)$ is $\cB$.} The positions $0$ and $3$ are $\not\equiv q_0\pmod 3$, so by (1)
$f(0),f(3)\in\{\cB,\cC\}$; they are not both $\cB$, by (AB) with difference $3$ ($h(3)=\cB$). By (2),
$f(H),f(H+3)\in\{\cB,\cC\}$. By (AB) with difference $H$ and $h(H)=\cC$ (Lemma~\ref{lem:tail}), $f(0),f(H)$ are not
both $\cC$, and $f(3),f(H+3)$ are not both $\cC$. If $f(0)=\cC$ then $f(H)=\cB$; otherwise $f(3)=\cC$ and
$f(H+3)=\cB$.

\emph{(4) $f(6)=\cC$.} By (1), $f(6)\ne\cA$. Suppose $f(6)=\cB$. Then $f(0)\ne\cB$ (difference $6$, $h(6)=\cB$) and
$f(3)\ne\cB$ (difference $3$), so $f(0)=f(3)=\cC$, and as in (3) $f(H)=f(H+3)=\cB$. But $f(H),f(H+3)$ are at
difference $3$ and $h(3)=\cB$: a contradiction.

\emph{(5) The colour of $L$.} Recall $c(L)=f(H+6)$.
\begin{itemize}
\item $f(H+6)\ne\cC$: (AB) with difference $H$ between positions $6$ and $H+6$, with $h(H)=\cC$ and $f(6)=\cC$.
\item $f(H+6)\ne\cB$: by (3) one of $f(H),f(H+3)$ is $\cB$, at difference $6$ or $3$ from $H+6$, and $h(6)=h(3)=\cB$.
\item $f(H+6)\ne\cA$: the difference $\delta=H+6-p_{\rm high}=a^2-2a+7-q_0$ satisfies $3a<\delta\le M$ (as $a\ge7$),
so $h(\delta)=h^\ast(\delta)=\cA$, while $f(p_{\rm high})=\cA$.
\end{itemize}
All three colours are excluded.
\end{proof}

%=====================================================================
\section{Proof of the main theorem}

\begin{proof}[Proof of Theorem~\ref{thm:upper}]
For $4\le a\le26$ this is Lemma~\ref{lem:small}. Let $a\ge27$ and suppose $c$ is a $3$-colouring of $[1,L]$ without
monochromatic solution. Its restriction to $[1,L-1]$ is solution-free, so by Corollary~\ref{cor:p0} its prefix is
$\Pz$. Section~\ref{sec:p0} (which needs $a\ge7$) then shows that $c(L)$ admits no colour, a contradiction.
\end{proof}

\begin{proof}[Proof of Theorem~\ref{thm:main}]
Theorem~\ref{thm:azb} gives $R_3(ax+3y=3z)\ge L$, and Theorem~\ref{thm:upper} gives $R_3(ax+3y=3z)\le L$.
\end{proof}

%=====================================================================
\section{Range conditions}\label{sec:ranges}

The following inequalities were used above. Each is elementary; all of them, and every BB-triple identity of
Sections~\ref{sec:53} and \ref{sec:p0}, are also checked mechanically for every $a\le400$ with $3\nmid a$ by
\texttt{tools/check\_hand\_proofs.py}.
\begin{enumerate}
\item $3H\le L-1$ and $1+aF=L-a$ (Section~\ref{sec:setup}).
\item For $a\ge23$: with $a=3m+1$ ($m\ge8$), $J=2m+4$, $m_0=6m^2+23m+17$, $F=9m^2+9m+8$ and $F-m_0-67=3m^2-14m-76\ge4$;
with $a=3m+2$ ($m\ge7$), $J=2m+4$, $m_0=6m^2+25m+22$, $F=9m^2+15m+12$ and $F-m_0-67=3m^2-10m-77\ge0$. Both
expressions increase with $m$. Also $m_0\le F-67<H$ and $147\le\min(H,F)$.
\item For $a\ge7$: $v_3=a-6\ge1$; $a^2+3\le M$; $(a-3)(a+3)+a\le M$; $3a<s_T\le M$; $a^2+a-12>3a$; $3a+3\le M$;
$(1+2a)/3+a^2\le M$; $a^2-2a-1>3a$; $3a<a^2-2a+7-q_0\le M$.
\end{enumerate}

%=====================================================================
\section{Certificates}\label{sec:cert}

\subsection{What is certified}
\begin{center}\begin{tabular}{llr}\toprule
claim & directory in \texttt{certificates/} & number\\\midrule
Lemma~\ref{lem:small}, $4\le a\le26$ & \texttt{small\_a/} & 16\\
Lemma~\ref{lem:W} & \texttt{lemma\_W/} & 2\\
Lemma~\ref{lem:r76}, late case ($\mathtt R_X$) & \texttt{third\_colour\_late/} & 53\\
Lemma~\ref{lem:r76}, early case ($\mathtt I_X$) & \texttt{third\_colour\_early/} & 53\\\midrule
total used in the proof & & 124\\
cross-check of Lemma~\ref{lem:W} & \texttt{lemma\_W\_alt/} & 2\\\bottomrule
\end{tabular}\end{center}

Every formula is in DIMACS format and comes with an LRAT proof of unsatisfiability produced by CaDiCaL~\cite{cadical}
(commit \texttt{c6073042}) with options \texttt{--lrat --no-binary}. Each proof was checked with \texttt{lrat-check}
from drat-trim (commit \texttt{2e3b2dc0}) and with the formally verified checker \texttt{cake\_lpr}~\cite{cakelpr}
(commit \texttt{a36874a8}). A proof checker only certifies that the given formula is unsatisfiable; that each formula
encodes what is stated above is ensured by the following.

\subsection{Faithfulness of the formulas}
\begin{itemize}
\item The formulas $\mathtt Ua$, $\mathtt I_X$, $\mathtt R_X$ are generated by short programs
(\texttt{small\_a/ucert.py} and \texttt{two64.py} in the two \texttt{third\_colour\_*} directories) that implement the definitions of Sections~\ref{sec:small}
and~\ref{sec:53} literally.
\item For Lemma~\ref{lem:W}, each formula comes with a table (\texttt{*.coords.json}) giving, for every variable, its
vertex as an integer polynomial in $k$. The program \texttt{lemma\_W/audit/audit\_free.py}, reading only the DIMACS
file and this table, checks that every clause is (i) an exactly-one-colour clause, (ii) an (AB) clause whose three
vertices satisfy $q'-q\equiv t$ as polynomials with $t\ge1$, (iii) a (BB) clause whose vertices satisfy
$z-y\equiv(3k+\rho)u$ as polynomials with $u\in\{1,2,k,k+1\}$, or (iv) one of the auxiliary clauses of the target;
and that every vertex lies in its range for every integer $k\ge5$. The program
\texttt{concrete\_free.py} instantiates every (AB)/(BB) clause for $5\le k\le60$ and checks on integers that it
comes from a solution of $ax+3y=3z$ inside $[1,L-1]$. The program \texttt{glue.py} enumerates all $3^8$
colourings of $h(1..8)$ and, for each, decides whether the target clause together with the clauses defining its
auxiliary variables is satisfiable. It checks that this is the case exactly for the colourings whose normalisation
is \emph{not} admissible; equivalently, the colourings that Lemma~\ref{lem:W} leaves open are exactly those whose
normalisation is one of the $53$ admissible words used in Lemma~\ref{lem:r76}.
\item The cross-check formulas in \texttt{lemma\_W\_alt/} contain the same base clauses as the two formulas in
\texttt{lemma\_W/} and express the target differently (one blocking clause per allowed colouring of $h(1..8)$); they are also
unsatisfiable.
\end{itemize}

\subsection{Replay and data availability}
All formulas, proofs, coordinate tables, generators and audit scripts are contained in the repository that
accompanies this document, together with a SHA-256 manifest of every file. The script
\texttt{certificates/replay.sh} checks the SHA-256 of every formula and runs both proof checkers on every
certificate; see \texttt{README.md}. On 18 October 2026 all $126$ certificates were replayed from a clean copy of the
repository with the checker versions above, with no failure.

%=====================================================================
\appendix
\section{The 53 admissible prefixes}\label{app:53}
For each admissible word $X$ the row length $K=K_X$ of Lemma~\ref{lem:r76} and the bound $F_X=K_X+79$ of the early
case.
\begin{center}
{\small
\begin{tabular}{lrr@{\qquad}lrr@{\qquad}lrr}\toprule
$X$ & $K$ & $F_X$ & $X$ & $K$ & $F_X$ & $X$ & $K$ & $F_X$\\\midrule
$\cA\cA\cA\cA\cA\cA\cB\cB$ & 68 & 147 & $\cA\cA\cB\cB\cA\cB\cA\cA$ & 17 & 96 & $\cA\cB\cA\cB\cA\cA\cB\cB$ & 15 & 94\\
$\cA\cA\cA\cA\cB\cA\cA\cB$ & 60 & 139 & $\cA\cA\cB\cB\cA\cB\cA\cB$ & 13 & 92 & $\cA\cB\cA\cB\cB\cA\cA\cA$ & 15 & 94\\
$\cA\cA\cA\cA\cB\cA\cB\cB$ & 25 & 104 & $\cA\cA\cB\cB\cA\cB\cB\cA$ & 14 & 93 & $\cA\cB\cA\cB\cB\cA\cA\cB$ & 15 & 94\\
$\cA\cA\cA\cB\cA\cA\cA\cB$ & 13 & 92 & $\cA\cA\cB\cB\cA\cB\cB\cB$ & 14 & 93 & $\cA\cB\cA\cB\cB\cA\cB\cA$ & 17 & 96\\
$\cA\cA\cA\cB\cA\cA\cB\cA$ & 51 & 130 & $\cA\cA\cB\cB\cB\cB\cA\cA$ & 17 & 96 & $\cA\cB\cA\cB\cB\cA\cB\cB$ & 17 & 96\\
$\cA\cA\cA\cB\cA\cA\cB\cB$ & 14 & 93 & $\cA\cA\cB\cB\cB\cB\cA\cB$ & 13 & 92 & $\cA\cB\cB\cA\cA\cB\cA\cA$ & 16 & 95\\
$\cA\cA\cA\cB\cB\cA\cA\cA$ & 41 & 120 & $\cA\cA\cB\cB\cB\cB\cB\cA$ & 14 & 93 & $\cA\cB\cB\cA\cA\cB\cA\cB$ & 15 & 94\\
$\cA\cA\cA\cB\cB\cA\cA\cB$ & 14 & 93 & $\cA\cA\cB\cB\cB\cB\cB\cB$ & 13 & 92 & $\cA\cB\cB\cA\cA\cB\cB\cA$ & 16 & 95\\
$\cA\cA\cA\cB\cB\cA\cB\cA$ & 21 & 100 & $\cA\cB\cA\cA\cA\cA\cA\cB$ & 27 & 106 & $\cA\cB\cB\cA\cA\cB\cB\cB$ & 15 & 94\\
$\cA\cA\cA\cB\cB\cA\cB\cB$ & 17 & 96 & $\cA\cB\cA\cA\cA\cA\cB\cA$ & 43 & 122 & $\cA\cB\cB\cA\cB\cB\cA\cA$ & 15 & 94\\
$\cA\cA\cB\cA\cA\cB\cA\cA$ & 10 & 89 & $\cA\cB\cA\cA\cA\cA\cB\cB$ & 22 & 101 & $\cA\cB\cB\cA\cB\cB\cA\cB$ & 20 & 99\\
$\cA\cA\cB\cA\cA\cB\cA\cB$ & 11 & 90 & $\cA\cB\cA\cA\cB\cA\cA\cA$ & 33 & 112 & $\cA\cB\cB\cA\cB\cB\cB\cA$ & 17 & 96\\
$\cA\cA\cB\cA\cA\cB\cB\cA$ & 13 & 92 & $\cA\cB\cA\cA\cB\cA\cA\cB$ & 25 & 104 & $\cA\cB\cB\cA\cB\cB\cB\cB$ & 24 & 103\\
$\cA\cA\cB\cA\cA\cB\cB\cB$ & 13 & 92 & $\cA\cB\cA\cA\cB\cA\cB\cA$ & 15 & 94 & $\cA\cB\cB\cB\cA\cB\cA\cA$ & 17 & 96\\
$\cA\cA\cB\cA\cB\cB\cA\cA$ & 17 & 96 & $\cA\cB\cA\cA\cB\cA\cB\cB$ & 16 & 95 & $\cA\cB\cB\cB\cA\cB\cB\cA$ & 17 & 96\\
$\cA\cA\cB\cA\cB\cB\cA\cB$ & 13 & 92 & $\cA\cB\cA\cB\cA\cA\cA\cA$ & 15 & 94 & $\cA\cB\cB\cB\cB\cB\cA\cA$ & 17 & 96\\
$\cA\cA\cB\cA\cB\cB\cB\cA$ & 14 & 93 & $\cA\cB\cA\cB\cA\cA\cA\cB$ & 15 & 94 & $\cA\cB\cB\cB\cB\cB\cB\cA$ & 44 & 123\\
$\cA\cA\cB\cA\cB\cB\cB\cB$ & 13 & 92 & $\cA\cB\cA\cB\cA\cA\cB\cA$ & 15 & 94 & & &\\
\bottomrule
\end{tabular}}
\end{center}
The eleven non-admissible words $\cA w_2\cdots w_8$ (with $w_3=w_6$) are those listed in Section~\ref{sec:W}.

\section*{Use of AI tools}
The proof strategy, the computer encodings, the checking scripts and the text of this document were developed with
substantial assistance from AI systems (Anthropic Claude and OpenAI GPT), which also cross-checked one another's
computations independently. All unsatisfiability claims rest on the certificates and checkers described in
Section~\ref{sec:cert}. This document has not been peer reviewed by a human expert.

\begin{thebibliography}{9}
\bibitem{AZB} T.~Ahmed, L.~Zaman, C.~Bright, \emph{Symbolic sets for proving bounds on Rado numbers}, in: Proc.\ SC$^2$
2025, CEUR Workshop Proceedings Vol.~4116 (2025); arXiv:2505.12085v3.
\bibitem{CDW} Y.~Chang, J.~A.~De~Loera, W.~J.~Wesley, \emph{Rado numbers and SAT computations}, in: ISSAC 2022,
ACM (2022), doi:10.1145/3476446.3535494; arXiv:2210.03262.
\bibitem{cadical} A.~Biere, T.~Faller, K.~Fazekas, M.~Fleury, N.~Froleyks, F.~Pollitt, \emph{CaDiCaL 2.0}, in: CAV 2024,
LNCS 14682, Springer (2024) 133--152.
\bibitem{LRAT} L.~Cruz-Filipe, M.~J.~H.~Heule, W.~A.~Hunt, M.~Kaufmann, P.~Schneider-Kamp, \emph{Efficient certified
RAT verification}, in: CADE-26, LNCS 10395, Springer (2017) 220--236.
\bibitem{cakelpr} Y.~K.~Tan, M.~J.~H.~Heule, M.~O.~Myreen, \emph{cake\_lpr: verified propagation redundancy checking in
CakeML}, in: TACAS 2021, LNCS 12652, Springer (2021) 223--241.
\end{thebibliography}
\end{document}
